\subsection{商空间与乘积空间中的半范数}

设 E 是 K 上的一个拓扑向量空间，其拓扑由半范数集 \(\Gamma\) 定义。显然，\(\Gamma\) 中诸半范数在 E 的向量子空间 M 上的限制，定义了由 E 的拓扑在 M 上诱导的拓扑。

设 \(\phi\) 是 E 到向量商空间 E/M 上的典范映射。我们证明，对 E 上的每个半范数 \(p\)，函数

\[
\dot {p} (z) = \inf _ {\phi (x) = z} p (x)\tag{3}
\]

是 E/M 上的半范数。事实上，显然 \(\dot{p}\) 满足条件 \((\mathrm{SN}_{\mathrm{I}})\)；另一方面，若 \(z'\)、\(z''\) 是 E/M 的两个向量，我们有：

\[
\begin{array}{r l} \inf _ {\phi (x) = z ^ {\prime} + z ^ {\prime \prime}} p (x) & \leqslant \inf _ {\phi (x ^ {\prime}) = z ^ {\prime}, \phi (x ^ {\prime \prime}) = z ^ {\prime \prime}} p (x ^ {\prime} + x ^ {\prime \prime}) \\ & \leqslant \inf _ {\phi (x ^ {\prime}) = z ^ {\prime}, \phi (x ^ {\prime \prime}) = z ^ {\prime \prime}} \left(p (x ^ {\prime}) + p (x ^ {\prime \prime})\right) \\ & = \inf _ {\phi (x ^ {\prime}) = z ^ {\prime}} p (x ^ {\prime}) + \inf _ {\phi (x ^ {\prime \prime}) = z ^ {\prime \prime}} p (x ^ {\prime \prime}) \end{array}
\]

这表明 \(\dot{p}\) 满足 \((\mathrm{SN}_{\mathrm{II}})\)。我们称 \(\dot{p}\) 是 p 关于 M 的商半范数。

同样的推理证明：若 p 是超半范数，则 \(\dot{p}\) 也是超半范数。

这样，我们有（采用第 2 目的记号）

\[
\phi (\mathrm{W} (p, \alpha)) = \mathrm{W} (\dot {p}, \alpha).\tag{4}
\]

对每个 \(\alpha > 0\) 成立。事实上，说 \(\dot{p}(z) < \alpha\)，意味着存在 \(x \in \mathbf{E}\) 使得 \(\phi(x) = z\) 且 \(p(x) < \alpha\)，由此即得关系式 (4)。

由此我们推出：若半范数集 \(\Gamma\) 是有向的（II, 第 3 页, 注记 3），则 E/M 上的商拓扑由诸半范数 \(\dot{p}\)（其中 p 在 \(\Gamma\) 中变化）所成的集定义。

若 N 是 \(\{0\}\) 在 E 中的闭包，则 E/N 的拓扑由诸商半范数 \(\dot{p}\)（其中 p 在 \(\Gamma\) 中变化）定义（即使 \(\Gamma\) 不是滤过的）：这里，\(\dot{p}(\dot{x}) = p(x)\)（对属于类 \(\dot{x} \mod N\) 的每个 x）。注意，E/N 不是别的，正是与 E 相伴的 Hausdorff 空间（I, 第 4 页）。

设 \(\mathbf{E}\) 是 \(\mathbf{K}\) 上的向量空间，\((\mathrm{E}_{\iota})_{\iota \in \mathbf{I}}\) 是 \(\mathbf{K}\) 上的一族向量空间，其中 \(\mathrm{E}_{\iota}\) 赋有拓扑 \(\mathcal{T}_{\iota}\)（由半范数集 \(\Gamma_{\iota}\) 定义的）。对每个 \(\iota \in \mathbf{I}\)，设 \(f_{\iota}\) 是 \(\mathbf{E}\) 到 \(\mathrm{E}_{\iota}\) 中的一个线性映射；显然，当 \(p_{\iota}\) 在集 \(\Gamma_{\iota}\) 中变化时，诸 \(p_{\iota} \circ f_{\iota}\) 构成一个半范数集 \(\Gamma_{\iota}'\)（\(\mathbf{E}\) 上的）。于是，拓扑 \(\mathcal{T}\)（\(\mathbf{E}\) 上的，定义为使所有映射 \(f_{\iota}\) 都连续的拓扑中最粗者；I, 第 9 页）由半范数集 \(\Gamma' = \bigcup_{\iota \in \mathbf{I}} \Gamma_{\iota}'\) 定义，这由 \(\mathcal{T}\) 的 0 的邻域的定义得出（GT, I, § 2.3, prop. 4）。

若诸 \(p_{\iota}\) 是超半范数，则诸 \(p_{\iota} \circ f_{\iota}\) 也是超半范数。

设 E 是 K 上的向量空间，赋有由半范数族 \((p_{\iota})_{\iota \in I}\) 定义的拓扑 T；对每个 \(\iota \in I\)，设 \(T_{\iota}\) 是由单个半范数 \(p_{\iota}\) 定义的拓扑，并用 \(E_{\iota}\) 表示由 E 赋有拓扑 \(T_{\iota}\) 而得的空间。则拓扑 T 是对角映射 \(\Delta : E \to \prod_{\iota \in I} E_{\iota}\) 下 \(\prod_{\iota \in I} \mathrm{E}_{\iota}\) 上乘积拓扑的原像（I, 第 9 页, prop. 7）。对每个 \(\iota \in I\)，用 \(\mathbf {N} _ {\iota}\) 表示 \(\{0 \}\) 在 \(\mathrm{E} _ {\iota}\) 中的闭包，并用 \(F_{\iota} = E_{\iota}/N_{\iota}\) 表示由范数 \(\hat{p}_{\iota}\)（对应于 \(p_{\iota}\) 的）定义的赋范空间（II, 第 4 页, 公式 (3)）；若 \(\phi_{\iota}: E_{\iota} \to F_{\iota}\) 是典范映射，\(\phi : (x_{\iota}) \mapsto (\phi_{\iota}(x_{\iota}))\) 是乘积映射，我们知道 \(\prod_{i \in I} E_{i}\) 上的乘积拓扑是 \(\phi\) 下 \(\prod_{\iota \in I}F_{\iota}\) 上乘积拓扑的原像（GT, II, § 3.9, prop. 18）。因此，拓扑 \(\mathcal{T}\) 是复合映射 \(\phi \circ \Delta\) 下 \(\prod_{\mathfrak{t}\in \mathbf{I}}\mathrm{F}_{\mathfrak{t}}\) 上乘积拓扑的原像。特别地，若 \(\mathcal{T}\) 是 Hausdorff 的，则由 II, 第 3 页, prop. 2 得 \(\Delta^{-1}(\prod_{\iota \in I} \mathbf{N}_{\iota}) = \{0 \}\)，因而映射 \(\phi \circ \Delta\) 是单射，因而：

命题 3. --- K 上每个其拓扑由半范数集定义的 Hausdorff 拓扑向量空间 E，都同构于一个 Banach 空间乘积的子空间。

进而，若 E 的拓扑由可数的半范数集定义，则 E 可度量化（I, 第 16 页）。

